% LaTeX-Vorlage Version 2.1 ‚

% TODO: individuelle Einstellungen (Name, Titel etc.)
% -> bitte in Konfigurationsdatei anpassen
\input{includes/config.tex}

% Grundlegende Dokumenteneigenschaften gemäß DHBW-Vorgaben
\documentclass[a4paper,fontsize=11pt,oneside,parskip=half,headings=normal,listof=nochaptergap]{scrreprt} 
% \usepackage{showframe} % nur für Kontrolle der Ränder 

%%% Präambel einbinden (mit Festlegungen gemäß DHBW-Vorgaben) %%%
\input{template/_dhbw_praeambel.tex}

%%% Name der eigenen Literatur-Datenbank (ggf. anpassen) %%%
\bibliography{includes/literatur-datenbank.bib}

\begin{document}
%%% Deckblatt gemäß DHBW-Vorgaben einbinden (keine Anpassung nötig) %%% 
\input{template/_dhbw_deckblatt.tex}

%%% Umstellung der Seiten-Nummerierung auf i, ii, iii ... %%%
\pagenumbering{Roman}

%%% Abstract einbinden (optionale Kurzfassung Ihrer Arbeit) %%%
% \input{includes/abstract.tex}
% \cleardoublepage

%%% Inhalts-, Abbildungs-, Tabellenverzeichnisse %%%
% werden einzeilig gesetzt, um Platz zu sparen 

\begin{spacing}{1}
    \tableofcontents % Inhaltsverzeichnis ausgeben
    \clearpage
    \startAbkVerzeichnis

\begin{acronym}[DSGVO] 
\acro{API}{Application Programming Interface}
\acro{DSGVO}{Datenschutz-Grundverordnung}
\acro{GeschGehG}{Geschäftsgeheimnisgesetz}
\acro{IT}{Informationstechnik}
\acro{KI}{Künstliche Intelligenz}
\acro{LLM}{Large Language Model}
\acro{NER}{Named Entity Recognition}
\acro{NIS-2}{Sicherheit von Netz- und Informationssystemen}
\acro{PII}{Personally Identifiable Information}
\acro{Regex}{Reguläre Ausdrücke}
\acro{REST}{Representational State Transfer}
\acro{SLA}{Service Level Agreement}
\acro{TLS}{Transport Layer Security} 
\end{acronym} % Abkürzungsverzeichnis einbinden

    % \clearpage
    % \thispagestyle{kapitelkopfzeile}
    % \listoffigures
    % \phantomsection
    % \addcontentsline{toc}{chapter}{\listfigurename} % Abb.verz. ins Inh.verz. aufnehmen

    \clearpage
    \listoftables
    \phantomsection
    \addcontentsline{toc}{chapter}{\listtablename} % Tab.verz. ins Inh.verz. aufnehmen
\end{spacing}

%%% Umstellung der Seiten-Nummerierung auf 1, 2, 3 ... %%%
\cleardoublepage
\pagenumbering{arabic}

%%% Kapitel ohne Seitenumbruch und kompaktere Überschriften %%%
\RedeclareSectionCommand[
  style=section, 
  indent=0pt,
  beforeskip=1ex plus .5ex minus .5ex,
  afterskip=0.5ex plus .1ex
]{chapter}
\RedeclareSectionCommands[
  beforeskip=1ex plus .5ex minus .5ex,
  afterskip=0.5ex plus .1ex
]{section,subsection,subsubsection}

%%% Ihr eigentlicher Inhalt %%%
% Empfehlung: strukturieren Sie Ihren Text in einzelnen Dateien 
% und binden Sie diese hier mit \input{includes/dateiname.tex} ein
%\chapter{Hauptteil}
\label{chapter:hauptteil}


\chapter{Einleitung}

\section{Ausgangslage und Problemstellung}
In den letzten Jahren hat der Umgang mit Künstliche Intelligenz (KI) zu einer fundamentalen Veränderung in der \ac{IT}-Landschaft geführt \footcite{naveed2023comprehensive}.
Gerade \ac{LLM}-Modelle haben als transformative Technologie die Art und Weise, wie wir arbeiten, wie wir Informationen beschaffen und wie wir Probleme lösen, revolutioniert \footcite{alzubaidi2024researchlandscape}.
Besonders beliebt ist dabei die Integration von bestehenden \ac{LLM}-\acp{API} in eigene Softwarelösungen, da Unternehmen so schnell und kostengünstig von den Vorteilen der \ac{KI} profitieren können \footcite{kucharavy2025impactllmassistants}.
So werden bei Bosch z. B. \ac{LLM}-\acp{API} verwendet, um die \ac{IT}-Infrastrukturen von Anwendungen zu analysieren und zu optimieren.

Das Problem ist dabei, dass die Integration von \ac{LLM}-\acp{API} in Softwarelösungen zwar starke Effizienzsteigerungen ermöglicht, gleichzeitig aber ein oft übersehenes Risiko mit sich bringt. Das Übertragen von kritischen Unternehmensdaten an Dritte korreliert unmittelbar mit unternehmenswirtschaftlichen Erfolgen. Wichtig dabei zu verstehen ist, dass die Verschlüsselung der Anfragen über \ac{TLS} zwar die Integrität und Vertraulichkeit der Daten sichert, nicht aber den Zugriff des Anbieters selbst \footcite{staab2025sokprivacyparadox}.

\section{Zielsetzung und Aufbau der Arbeit}
Um dieses Risiko zu minimieren, gilt es als Aufgabe dieser Arbeit, eine unsichtbare Schicht zwischen der Software und der \ac{API} zu schaffen, mit der kritische Daten maskiert werden. So kann gewährleistet werden, dass die Daten dem \ac{API}-Provider nicht ungeschützt zugänglich sind.

Die Frage, die diese Arbeit dabei beantworten soll, lautet dabei:
\begin{equation}
    \parbox{0.85\textwidth}{\centering \textit{Inwiefern ermöglicht ein hybrider Data-Masking-Ansatz (\ac{Regex} + \ac{NER}) einen effektiven Schutz sensibler \ac{IT}-Infrastrukturdaten bei gleichzeitig hoher Antwortqualität von \ac{LLM}-\acp{API}?}}
\end{equation}
Um dieser nachzugehen, werden zunächst die theoretischen Grundlagen der \ac{LLM}-\ac{API}-Kommunikation und die damit verbundenen Risiken analysiert. Anschließend werden verschiedene Maskierungsansätze untersucht und unter anderem deren Einfluss und Effektivität auf die Antwortqualität untersucht \footcite{zhang2025surveyprivacypreservation}. Dabei schließt die Arbeit mit einer Bewertung der praktischen Einsetzbarkeit unter Berücksichtigung regulatorischer und governance-relevanter Aspekte.
\chapter{Theoretische Grundlagen und Problemanalyse}

\section{Funktionsweise von \ac{LLM}-\acp{API}}

Die Kommunikation mit \ac{LLM}-\acp{API} erfolgt typischerweise über \ac{REST}ful \acp{API}, bei denen Anfragen
an zentrale Server der \ac{API}-Provider gesendet werden.
Die Anfragen enthalten dabei den sogenannten Prompt sowie verschiedene Parameter wie Temperatur, Maximalzahl an Tokens, usw.
Die Verbindung zum \ac{API}-Provider wird typischerweise mit \ac{TLS} verschlüsselt und verhindert somit
Man-in-the-Middle-Angriffe und das Abhören von Dritten.
Da jedoch unverschlüsselte Daten für die Verarbeitung der Modelle notwendig sind,
kann der Zugriff des \ac{API}-Providers auf die Informationen nicht verhindert werden (European Union Agency for Cybersecurity, 2022).
Trotz \acp{SLA}s, welche die Sicherheit der Daten garantieren sollen,
entsteht ein Vertrauensproblem zwischen \ac{API}-Provider und \ac{API}-Konsumierer,
da die Kontrolle über die Daten bei externen Parteien liegt (Kroll \& Zittrain, 2023). Dies stellt vor allem ein Problem dar,
wenn die \ac{API}-Provider unter einer anderen regulatorischen oder wettbewerbswirksamen Umgebung als der \ac{API}-Konsumierer stehen.


\section{Spezifische Risiken in der \ac{IT}-Infrastruktur}
Um die Bedeutung der Daten zu verstehen, lohnt es sich diese zu kategorisieren.
Es wird unterschieden zwischen personenbezogenen Daten (\ac{PII}) einerseits und
geschäftskritischen Infrastrukturdaten andererseits. Personenbezogene Daten umfassen dabei alle Daten, die eine natürliche Person eindeutig identifizieren können.
Das können zum Beispiel Namen, E-Mail-Adressen, Telefonnummern, etc. sein. Bei diesen greift die \ac{DSGVO},
weswegen eine unachtsame Verwendung der Daten zu erheblichen Geldbußen und Reputationsschäden führen kann.
Die Kategorie der geschäftskritischen Infrastrukturdaten umfasst hingegen alle Daten, die für die Betriebsführung relevant sind.
Dazu zählen in diesem Fall unter anderem IP-Adressen und Netzwerksegmentierungen, Servernamen und deren Betriebssystemversionen, etc.
Im Gegensatz zu personenbezogenen Daten existieren dabei keine regulatorischen Rahmenbedingungen, die die Verwendung der Daten einschränken.
Das Risiko liegt stattdessen darin, dass die Offenlegung dieser die Sicherheitsinfrastruktur des Unternehmens massiv
kompromittieren kann (Microsoft Research, 2023).
\chapter{Technische Implementierung von Data-Masking}

\section{Verfahren der Anonymisierung}

Vor der detaillierten Betrachtung der einzelnen Verfahren bietet sich ein systematischer Vergleich der beiden zentralen Maskierungstechniken an:

\begin{table}[ht]
    \centering
    \footnotesize
    \caption{Vergleich der Maskierungstechniken}
    \label{tab:maskierungstechniken}
    \begin{tabular}{@{} >{\raggedright\arraybackslash}p{3.5cm} >{\raggedright\arraybackslash}p{5.5cm} >{\raggedright\arraybackslash}p{5.5cm} @{}}
        \toprule
        \textbf{Kriterium}    & \textbf{\ac{Regex}}                                 & \textbf{\ac{NER}}                            \\
        \midrule
        Erkennungsgenauigkeit & Hoch für strukturierte Daten (IP-Adressen, E-Mails) & Hoch für unstrukturierte, kontextuelle Daten \\
        Latenz                & Sehr gering (Millisekunden)                         & Höher (abhängig von Modellgröße)             \\
        Komplexität           & Niedrig (regelbasiert)                              & Hoch (KI-basiert, erfordert Training)        \\
        Wartbarkeit           & Manuelle Regelpflege erforderlich                   & Kontinuierliches Fine-Tuning nötig           \\
        Einsatzgebiet         & Formale, standardisierte Datenformate               & Freitext, variable Formulierungen            \\
        Ressourcenbedarf      & Minimal (CPU-basiert)                               & Hoch (GPU für Transformer-Modelle)           \\
        \bottomrule
    \end{tabular}
\end{table}

\subsection{Reguläre Ausdrücke}
\ac{Regex} gilt als verbreitetste Methode zur Identifikation von Mustern in Texten.
Es verwendet vordefinierte Zeichenketten, um spezifische Textmuster mit hoher Präzision in einer Zeichenkette zu finden.
Im Kontext von IT-Infrastrukturen können so z. B. IP-Adressen identifiziert werden und für die weitere Verarbeitung
mit anonymisierten Werten maskiert werden.

Der große Vorteil von \ac{Regex} ist dabei die schnelle und ressourcenschonende Verarbeitung der Daten, da moderne Regex-Engines hochoptimiert
für String-Verarbeitungsalgorithmen sind. Dabei können auch Datenbank-Verbindungsstrings, welche typischerweise Schlüsselwörter wie "server",
"database" oder "password" enthalten, zuverlässig erkannt und nachfolgende Informationen maskiert werden.
Die Implementierung einer solchen Lösung erfolgt typischerweise durch die Verwendung von \ac{Regex}-Bibliotheken oder externen Anbietern.

Die Limitation von \ac{Regex} liegt jedoch darin, dass es sich um eine rein regelbasierte Methode handelt.
So können z. B. Servernamen, welche einer bestimmten Namenskonvention folgen (z. B. "prod-db-server-01"), zuverlässig erkannt und maskiert werden.
Werden die Servernamen jedoch willkürlich gewählt (z. B. "der Hauptdatenbankserver"), kann \ac{Regex} diese nicht mehr erkennen,
da in diesem Fall semantische Zusammenhänge verwendet werden müssen und syntaktische Mustererkennung nicht ausreicht (Liu \& Vaidya, 2022).
Gleichzeitig besteht die Gefahr des Undermaskings: Neue oder unbekannte Namenskonventionen oder Datenformaten
werden nicht in die \ac{Regex}-Regeln aufgenommen und sommit nicht erkannt,
bzw. die Gefahr des Overmaskings, dass zu viele Daten maskiert werden und folglich die Antwortqualität der \ac{LLM}-\ac{API} sinkt.


\subsection{Named Entity Recognition}
Im Gegensatz zu \ac{Regex} adressiert \ac{NER} das Problem der Datenmaskierung mithilfe von semantischer Klassifikation.
Dabei wird ein kontextuelles Sprachverständnis genutzt, um die Bedeutung von Wörtern im umgebenden Text zu erfassen und somit auch
Entitäten zu erkennen, welche nicht explizit in Regeln definiert sind. \ac{NER} "denkt mit" und ein Servername wird nicht erkannt,
weil er aussieht wie ein Servername, sondern weil der Satzbau („...verbindet sich mit dem Host Alpha-Prime...“) darauf hindeutet.

Der Nachteil von \ac{NER} zeigt sich jedoch in der Komplexität und den Ressourcen, welche für das Training und den Betrieb benötigt werden.
So müssen die Modelle aufwendig trainiert und gewartet werden, was zeit- und kostenintensiv sein kann.
Aus diesem Grund lohnt sich der Einsatz von \ac{NER} nur, wenn eine hohe Genauigkeit bei der Erkennung von Entitäten erforderlich ist und die Ressourcen für das Training und den Betrieb vorhanden sind.

Als Alternative gibt es einen hybriden Kaskaden-Ansatz. Dabei wird die Anfrage zuerst mit minimaler Latenz
durch ein \ac{Regex}-basiertes Modell geschickt. Alles, was harte Muster enthält, wird erkannt und mit anonymisierten Namenskonventionen ersetzt.
Der verbleibende Text wird daraufhin an das \ac{NER}-basierte Sicherheitsnetz geschickt, welches verbleibende Informationen erkennt und maskiert.
Dieses System stellt dabei nicht nur eine hohe Effizienz bereit, sondern auch eine erhöhte Zuverlässigkeit, da sich die beiden Modelle ergänzen.
Um die anonymisierten Antworten der \ac{LLM}-\ac{API} wieder ersetzen zu können, ist es wichtig, einen gemeinsamen Token-Speicher zu nutzen, welcher
alle anonymisierten Begriffe speichert.


\section{Das Utility-Privacy-Dilemma}

Die zentrale Herausforderung bei einer erfolgreichen Datenmaskierung liegt in der Balance zwischen dem Schutz sensibler Daten (Privacy) und der
gleichzeitigen Erhaltung der Nützlichkeit der Daten (Utility) für die \ac{LLM}-\ac{API}. In der Wissenschaft
wird dies oft als Privacy-Utility-Trade-off (Dwork \& Roth, 2014) bezeichnet.
Maximale Antwortqualität für die \ac{LLM}-\ac{API} bedeutet maximale Datenmenge,
was wiederum maximale Risiken für die Privatsphäre bedeutet.
Umgekehrt bedeutet maximale Datenminimierung maximale Ungenauigkeit bei der Antwortqualität der \ac{LLM}-\ac{API}. So kann die \ac{LLM}-\ac{API} zum Beispiel
nicht mehr die Zusammenhänge der \ac{IT}-Infrastruktur verstehen und so fehlerhafte Antworten zurückgeben.
Als Lösung wird häufig ein „Label-Preserving Masking“ eingesetzt. Anstatt sensible Daten durch anonymisierte Werte (z. B. [Mask]) zu ersetzen,
werden diese durch Label (z.B. [IP-ADDRESS1]) ersetzt, welche im Kontext der Anfrage verstanden werden können. So bleibt die logische
Topologie der \ac{IT}-Infrastruktur erhalten und die Antwortqualität der \ac{LLM}-\ac{API} nicht beeinträchtigt.

\chapter{Kritische Bewertung und Praxisrelevanz}

\section{Governance-Aspekte: Einbindung in die \ac{IT}-Infrastruktur}
Wichtig für eine erfolgreiche Einbindung eines Maskierungsalgorithmus ist die Vermeidung von sogenannten Schatten-\acp{KI}. Darunter versteht man \ac{LLM}-Systeme, welche nicht von einer \ac{IT}-Abteilung bereitgestellt wurden und somit auch keinen Maskierungsalgorithmus verwenden \footcite{mansner2025shadowai}.
Da Mitarbeiter bei fehlenden oder unperformanten Systemen oft stattdessen auf nicht geschützte \ac{LLM}-Systeme zurückgreifen, ist es wichtig, interne \ac{LLM}-Systeme bereitzustellen, die nicht nur aktiv beworben werden, sondern auch auf dem aktuellen Stand der Technik sind.

Der hybride Ansatz ist hierbei entscheidend, da er die hohe Präzision von \ac{NER} mit der minimalen Latenz von \ac{Regex} kombiniert. Nur eine performante Lösung, die den Arbeitsfluss nicht durch lange Wartezeiten unterbricht, kann die Akzeptanz der Mitarbeiter sicherstellen und somit den Rückzug in ungeschützte Schatten-KI-Systeme effektiv verhindern \footcite{naseer2025recap}.

Zudem sollte man abwägen, ob ein lokales \ac{LLM}-Modell (on-premise), zum Beispiel auf eigenen Servern, günstiger und effizienter wäre, als ein pseudonymisierendes externes \ac{LLM}-Modell, welches über eine \ac{API} angesprochen wird \footcite{costbenefit2025onprem}. Dieses Modell würde keinen Maskierungsalgorithmus verwenden, was die Antwortqualität und Latenz verbessern würde.

Zuletzt ist es auch wichtig, dass Mitarbeiter über „Residual Risks“ (Restrisiken) geschult werden, denn kein Maskierungsalgorithmus kann die absolute Sicherheit der Daten garantieren \footcite{staab2025sokprivacyparadox}. Deshalb ist es wichtig, nicht nur die technische Implementierung, sondern auch die organisatorischen Maßnahmen zu prüfen.

[Image of cloud vs on-premise LLM infrastructure comparison]

\section{Compliance: \ac{DSGVO} \& Privacy by Design}
Auch die \ac{DSGVO} nennt explizit die datenschutzfreundliche Voreinstellung und den Datenschutz durch Technik (Privacy by Design) als Grundsätze der Datenverarbeitung. In Artikel 25 der \ac{DSGVO} heißt es dazu:
\begin{quote}
  Der Verantwortliche trifft geeignete technische und organisatorische Maßnahmen, um sicherzustellen, dass grundsätzlich nur Datenverarbeitungen vorgenommen werden, die für den jeweiligen Zweck erforderlich sind.
\end{quote}
Gleichzeitig wird in Artikel 32 der \ac{DSGVO} sogar explizit die Pseudonymisierung als Maßnahme genannt, welche zur Erfüllung dieser Pflicht beitragen können.
Zusätzlich wird in dem neuen EU AI Act die Verwendung von \ac{KI}-Systemen zur Analyse von biometrischen Daten oder zur Erkennung von Emotionen stark eingeschränkt, was zeigt, dass die Bedeutung von Datenschutz und Datensicherheit in der EU zunimmt \footcite{europeanparliament2025interplayaact}.
\chapter{Fazit und Handlungsempfehlung}

\section{Zusammenfassung der Ergebnisse}
Um die Kontrolle über interne Daten zu behalten und gleichzeitig den steigenden Compliance-Anforderungen gerecht zu werden, ist eine technische Pseudonymisierung bei der Nutzung externer \ac{LLM}-\acp{API} für Unternehmen heute unerlässlich \footcite{staab2025sokprivacyparadox}. Mit einem hybriden Ansatz, der die deterministische Präzision von \ac{Regex} für strukturierte Daten mit der semantischen Flexibilität von \ac{NER} für unstrukturierte Kontexte kombiniert, können die Vorteile beider Ansätze genutzt werden \footcite{naseer2025recap}.Trotzdem sollte beachtet werden,
dass eine vollständige Pseudonymisierung oft schwierig ist und die
Qualität der Ergebnisse beeinträchtigen kann. Deshalb sind kontinuierliche Optimierungen notwendig, um ein zu striktes Maskierungsverfahren zu vermeiden \footcite{wang2025revisitingtradeoffs}.

\section{Handlungsempfehlung für die Praxis}
Um ein solches Modell in der Praxis umzusetzen, gibt es verschiedene Anbieter, welche Open-Source-Lösungen anbieten. Zu nennen sind hierbei insbesondere Microsoft Presidio, Skyflow oder Privacera \footcite{vignati2024towardsmiddleware}. Dabei kann zum Beispiel ein Stufenmodell eingesetzt werden, welches eine von Grund auf strikte Maskierung sukzessiv je nach Anwendungsfall lockert. Wichtig dabei sind vor allem regelmäßige Audits der Maskierungs-Logs, um Fehlklassifizierungen der Modelle zu finden \footcite{zhang2025surveyprivacypreservation}.
%\input{includes/text_mit_zitaten.tex}
%\input{includes/abbildungen_und_tabellen.tex}
%\input{includes/anhang.tex}
%\input{includes/release_notes.tex}
%%% Ende des eigentlichen Inhalts %%%


%%% Quellenverzeichnisse (keine Anpassung nötig) %%%
\clearpage
\literaturverzeichnis
%%% Ende Quellenverzeichnisse %%%


%%% Erklärung (keine Anpassungen nötig) %%%
% steht ganz am Ende des Dokuments
\cleardoublepage
\input{includes/erklaerung_ki.tex}
\input{template/_dhbw_erklaerung.tex}
\end{document}
